\documentclass[11pt]{article}
\usepackage{booktabs,amsmath,amssymb,graphicx}
\usepackage[margin=1in]{geometry}
\usepackage[hidelinks]{hyperref}
\title{PSI: Position-Stable Indexing for Long-Context Sparse Attention\\
\large with Upper-Bound Block Screening and Far/Near Budget Partitioning}
\date{2026-10-05}
\begin{document}
\maketitle

% =====================================================================
% Abstract: follows the MoBA four-sentence pattern (stakes + obstacle ->
% critique of two existing method families -> we propose + key mechanism
% + headline numbers -> deployment statement).
% =====================================================================
\begin{abstract}
Extending the effective context length of large language models is bounded by KV-cache bandwidth: with a 131K context on an 8B model, the full KV read per decode step dominates compute and caps throughput. Existing solutions face a hard trade-off. Training-free sparse indexers rely on lower-bound score approximations and collapse on long-context retrieval (we measure Quest at 21 on RULER 16K multikey\_3); training-based indexers such as DSA introduce learnable parameters and a thousand-step warm-up, and their selection quality depends on the training distribution. This paper proposes PSI, a training-free two-level sparse attention indexer. The first level scores 64-token blocks with a minmax \textbf{upper bound} over a 4bit-quantized position-stable subspace, which guarantees no miss within the subspace-score protocol; the second level partitions the token budget into a far pool and a near pool so that distant retrieval-critical tokens survive competition with near high scorers. On 13 LongBench tasks with Qwen3-8B, PSI is on par with FullKV in accuracy (50.78, with the 95\% per-sample paired-bootstrap CI containing 0); on the far-heavy multi-hop task musique it reaches 34.71, the highest score among all methods including FullKV (the single-task CI also contains 0). On RULER the synthetic-needle reference arm reaches 87.93 and the main arm 85.83, with the gap concentrated on far-dense common-word tasks and stable across lengths; the partition gain turns positive at 16K ($+0.91$), and both arms are reported by task profile. On the speed side, under this accuracy parity the selection chain achieves a kernel-level speedup of 1.46$\times$ at 32K to 5.09$\times$ at 128K over dense scoring; the e2e three-arm comparison (same machine, same protocol, 30B-A3B) is reported as measured: Quest wins e2e in every shape (TP2 64K total 62.1s versus 142.0s for PSI), and the only shape where PSI holds an advantage over Quest is the decode-segment index overhead (1.85$\times$); the speed claim is accordingly narrowed to the kernel, compute, and storage layers, with prefill index construction as the current boundary of any net e2e win. All accuracy costs and negative results are reported under matched hardware and matched protocol. PSI ships as an integration in sglang in production form (paged addressing, incremental indexing, CUDA graph).
\end{abstract}

% =====================================================================
% 1 Introduction: follows the MoBA six-paragraph pattern.
% P1 stakes + application anchors -> P2 technical challenge + direction
% -> P3 critique of three existing families -> P4 research question
% -> P5 we introduce + mechanism -> P6 contribution list.
% =====================================================================
\section{Introduction}

% P1 (stakes): long context is a core capability; KV bandwidth is the structural bottleneck.
The ability to process long context is becoming a core competitive capability of large language models. Long-chain reasoning in DeepSeek-R1~\cite{dsr1}, million-token context products from Kimi and Claude, and long-document question answering and repository-scale code understanding all require a model to consume tens of thousands to hundreds of thousands of tokens within a single request. Every extension of the context directly inflates the KV cache: with a 131K context on an 8B model, the KV read already dominates per-step compute, so the full KV access in decode forms a bandwidth bottleneck, and prefill is likewise constrained by full-attention bandwidth. Inference-system advances (PagedAttention~\cite{vllm}, FlashAttention~\cite{fa}) optimize memory organization and compute kernels, but they do not reduce the amount of data that must be read.

% P2 (challenge + direction): attention sparsity is the way to cut data volume; the three-layer mass distribution is our empirical starting point.
The way to cut KV reads is attention sparsity: softmax normalization concentrates the overwhelming majority of attention mass on a small number of tokens. Our trace analysis on real data with Qwen3-8B gives the concrete shape of this distribution: the mass of the dense top-1024 selection decomposes into three layers, sink (0.37--0.71), a near sliding window (0.16--0.29), and distant retrieval (0.02--0.29), and the distant layer profile correlates across layers within a task at 0.93--1.00. This three-layer structure implies the following: a sparse indexer that preserves quality must simultaneously satisfy three kinds of attention with completely different properties, and their budget demands conflict with each other.

% P3 (critique of three families): each family gets mechanism + quantified deficiency.
Existing solutions fall short along three paths. The first family, training-free sparse indexing, is represented by Quest~\cite{quest}, which constructs a score \textbf{lower bound} from per-page sampled minima. A lower bound underestimates genuinely high-scoring pages, and miss rates worsen as the number of distractor keys grows superlinearly with length; we measure this family collapsing to 21 on RULER 16K multikey\_3 (FullKV scores 100). KV compression methods (SnapKV, H2O, PyramidKV~\cite{snapkv,h2o,pyramidkv}) discard KV once at prefill and cannot respond to the decode-side query distribution. The second family, training-based indexing, is represented by DSA~\cite{dsa} (DeepSeek V3.2), which trades a learnable projection and ReLU scoring for selection capability at the cost of a 1000-step warm-up and a 2.1B-token training budget, with selection quality tied to the training distribution. The third family, the two-level structure (HISA~\cite{hisa}, COLM 2026), validates the block-coarse-plus-token-refine form, but two orthogonal design dimensions within it, whether the score representation takes an upper bound or a lower bound, and whether the budget is organized as a partition or a single pool, have not been systematically characterized for their quality and speed impact.

% P4 (research question, MoBA style): the critical research question.
This raises a critical research question: can we design a training-free sparse attention indexer with zero warm-up that, without introducing any learnable parameter, simultaneously preserves long-context retrieval quality and delivers the bandwidth win?

% P5 (we introduce + full mechanism paragraph).
This paper proposes PSI (Position-Stable Indexing) to answer this question. PSI is a two-level cascaded indexer. The first level organizes the KV cache into 64-token blocks, maintains per-block min/max bounds over a position-stable low-frequency subspace ($d'{=}32$) stored in 4bit quantized form, and scores each query against the optimistic side of the block bound: the upper bound guarantees no miss, so the block containing any genuinely high-scoring token is always admitted. The second level splits the token budget into a far pool and a near pool: the near sliding window is force-covered for locality, and the far pool applies token-level fine selection to the blocks admitted by L1, with an independent quota so that distant retrieval-critical tokens are not evicted by near high scorers. A per-request perception gate closes the chain: one softmax statistic from the last prefill segment identifies layers whose far mass is near zero, and decode safely skips their far retrieval. The entire selection chain is kernelized as a fused L1 block scorer, an L2 cascaded dual topk, and a unified sparse attention kernel, and an O(n) exact incremental index makes block-bound maintenance independent of sequence length.

% P6 (contribution list, claim-first).
Our contributions are as follows:
\begin{enumerate}
\item \textbf{A systematic characterization of training-free dimensionality-reduction freedom.} The position-stable low-frequency subspace is lossless for indexing quality (selection coverage $C{=}0.729$ over the full-dimensional top-1024, versus 0.732 for the same pipeline at $d'{=}128$ (4bit quantization and cascade losses included), against 0.32 for a random subspace). A triple verdict, measured in far-region mass recall ($R_{\mathrm{far}}$), confirms that the power of tail32 comes from \textbf{the 16 complete lowest-frequency RoPE rotation pairs}: frequency monotonicity ($R_{\mathrm{far}}$: low-frequency pairs 0.811, mid-frequency 0.360, high-frequency 0.004), pairing completeness (a misaligned pairing collapses to 0.567), and saturation at 16 pairs (8 pairs reach only 0.473). A four-quadrant characterization of selection methods bounds the feasible region: per-dimension spectral selection (0.759) and trained projections both lose to pairing entanglement; query-adaptive pair selection (0.865) cannot be indexed; per-layer static pairs (0.849) are index-compatible yet reverse at e2e ($+3.8$ points in replay, $-0.50$ at e2e), so tail32 is the e2e-verified zero-signal-cost optimum. Unsupervised PCA compresses fine selection further to $d{=}16$ at near-loss ($R_{\mathrm{far}}$ 0.766 versus same-batch tail32 0.804), and a shared SVD basis across kv-heads matches tail32 on 8/8 samples. Trained projections hold up in trace replay (92--94\%) yet collapse end-to-end: the projection basis is fragile under decode query-distribution drift, whereas the rotation-pair frequency prior is independent of the query distribution and stays stable. Dimensionality-reduction freedom comes from position priors, not from learning or adaptation.
\item \textbf{Far/near budget partition.} The selection index is built per kv-head: query heads in the same group share one kv-head-level selection, with no aggregation needed. Under GQA the distant attention mass is extremely uneven across kv-heads (a single kv-head can hold 0.804 of the far attention mass), so single-pool top-k evicts distant retrieval-critical tokens in favor of near high scorers. The partitioned independent quota holds on both e2e benchmarks (LongBench +0.62 = 50.78$-$50.16 and RULER +0.17 = 87.93$-$87.76 over the single pool, both directly verifiable from the main tables; these are internal ablation-arm comparisons with large task-level variance, read as structural signals rather than statistical significance), with the purest gain on far-heavy multi-hop tasks (musique 34.71 versus 29.74 for the same-budget single-pool arm, $+4.97$). Against FullKV the claim is accuracy parity at the same budget: the $+0.42$ 13-task difference has a 95\% per-sample paired-bootstrap CI of $[-0.18, +1.02]$ containing 0 (sign test 8/13, $p{=}0.58$), under which the selection chain still delivers a kernel speedup of 1.46$\times$ at 32K to 5.09$\times$ at 128K over dense scoring. Cluster-representative selection shows no consistent advantage under a strict token budget and is demoted to a negative-result ablation. An ablation of the within-group sharing itself (E103): per-q-head independent selection at the same $K_2$ loses accuracy instead (hotpotqa $-1.65$, musique $-0.20$ -- the 4 q-heads of a GQA group are semantically close, so the within-group mean aggregation acts as a 4-sample denoiser of the block/token importance estimate, and independent selection adds selection noise rather than information), so sharing saves a factor of 4 in index size while improving quality.
\item \textbf{A per-request perception gate.} Static layer masks do not transfer across tasks (multi-hop tasks lose 4.8--5.9 points); a dynamic signal from the last prefill segment correlates with decode-side far quality at 0.924, and in the real triggering regime ($\tau{=}0.005$), the gate-on versus gate-off difference is $+0.23$ on musique and $-0.41$ on qasper, and no far-heavy task collapses. The gate is positioned as a conservative switch that saves compute safely, not as a source of accuracy gain; the conservative arm ($\tau{=}0.05$) skips 14.3\% of layers at a far-mass loss of 0.8\% and saves selection-chain compute in equal proportion (upper bound: 13/36 layers from the offline profile).
\item \textbf{Tuning-free lightweight configuration selection.} The quality surface of the partition parameters $\alpha/\beta/\gamma$ is flat at four granularities -- task, layer, sample, and configuration arm: the per-task oracle of choosing them freely gains only $+0.08$ over the fixed default (13 tasks), the per-layer oracle only $+0.0008$, per-request online dispatch driven by a one-shot prefill far-statistic nets zero, and the per-arm oracle over 13 end-to-end configuration arms gains only $+0.21$ to $+0.49$; moreover, the rank correlation between the offline quality proxy and end-to-end accuracy is only 0.166 (confidence interval containing 0), and picking arms by the proxy loses 1.18 against the end-to-end optimum -- flatness is a structural property of the parameter surface, not a measurement artifact. Deployment therefore needs one default configuration within each task profile, with at most a two-point $\beta/\gamma$ selection between profiles (the real-task profile versus the synthetic-needle profile) -- still tuning-free in magnitude against the 1000-step warm-up of DSA, and the absence of tuning is itself a deployment advantage.
\end{enumerate}

% =====================================================================
% 2 Observations: three positional structure laws, each in the
% four-element pattern (phenomenon -> mechanism -> design implication
% -> boundary).
% =====================================================================
\section{Observations: Three Positional Structure Laws}\label{sec:obs}

The three designs of PSI each rest on a reproducible observation. Every observation is developed in four elements: phenomenon (trace statistics on real models and real tasks), mechanism (a structural explanation), design implication (what it directly dictates in PSI), and boundary (when it fails). The observations come from 10 real traces of Qwen3-8B (8 LongBench tasks plus needle/natural synthetic controls, saving all K and the trailing 256 queries per layer), with key laws re-verified on Qwen3-32B. Systematic negative results are part of our observation methodology: most papers report only when a law holds; we report also when it fails, and the four boundaries (protocol gap, trained-projection e2e collapse, adaptation having no leverage, nonlinear reduction being infeasible) come from the same protocol as the positive laws.

\subsection{Observation 1: Position-Segmented Density Steps of Attention Mass}\label{sec:obsa}

The mass captured by the dense top-1024 selection forms three positional steps (Figure~\ref{fig:h1}): sink 0.37--0.71, near sliding window 0.16--0.29, and far retrieval 0.02--0.29. Inside the far segment the mass is highly sparse and extremely uneven across GQA kv-heads: in layer 3 a single kv-head holds 0.804 of the far mass while the others in its group are near zero. The far-layer profile correlates 0.93--1.00 across layers within a task but only 0.05--0.89 across tasks: far retrieval is a task fingerprint, shared across layers, differentiated across tasks.

The two segments have opposite statistics, hence opposite optimal indexing granularity. The near segment is dense and stationary: candidate scores are uniformly high and close to each other, so a coarse index (block averages) misses nothing; its quality is flat from $d{=}4$ to $d{=}64$ under compression, so information is redundant and discriminability low. The far segment is sparse and sharp: a few critical tokens (the multi-hop evidence sentences of musique) carry almost all far mass, and missing one is expensive, which demands a no-miss block upper bound plus token-level fine selection. Sink and sliding window are the scoring-free extreme of the same spectrum: large mass at known positions. The partition is therefore the natural mapping of the density steps, not an arbitrary three-way cut.

Design implication: the far region takes minmax upper-bound coarse selection plus token-level fine selection, and the near region takes avg block scores that halve index storage (\S\ref{sec:l2}). The boundary must be stated honestly: the heterogeneous method combination itself contributes little (near-side method freedom $\leq 0.003$ in 16-sample wide-budget replay); the real gain of the partition is budget anti-eviction. Universality: the structure replicates on both the 8B and 32B models across 13 real tasks and synthetic needle (32B reproduces the far profile, with static layer skipping at 41/64 layers and precision 0.993), while the density ratios are task-dependent.

\begin{figure}[h]\centering
\includegraphics[width=0.95\linewidth]{figures/fig1_h1_decomposition.pdf}
\caption{Position-segmented density steps of attention mass (10 real traces): sink / near sliding window / far retrieval form the mass of the dense top-1024 selection; the far segment is sparse and extremely uneven across kv-heads.}
\label{fig:h1}
\end{figure}

\subsection{Observation 2: Budget Saturation and Flatness}\label{sec:obsb}

Three budget laws hold simultaneously (Figure~\ref{fig:abg}). First, far budget saturates: mass capture saturates after a far quota of 128--256 tokens; more budget buys no mass. Second, concession pays: a near budget share above the near-region share ($\beta{=}0.375$ over 0.25) holds across all real-task families, with the most far-heavy task musique gaining the most (+1.94). Third, the parameter surface is flat: $\beta$ is flat over 0.25--0.75, $\alpha$ is monotone toward its boundary, and the oracle upper bound of per-task perfect parameter selection is only $+0.08$, with $+0.0008$ per layer and $\le +0.49$ per configuration arm.

Together these laws characterize a structural waste of single-pool top-k methods. In the pooled competition of Quest and the single-pool form of TIA, the abundant mid-to-high scorers of the near segment occupy the quota, and the few critical far tokens at 0.02--0.29 are systematically evicted for their lower absolute scores. Once the far budget is independent it saturates at 128--256. The concession space is squeezed from both ends: far does not need that much (saturation), and near crowds out the rest (eviction), so budget should concede toward near. The positive translation of flatness is deployment simplicity: the quadruple negation of the oracles bounds perfect-versus-fixed configuration gaps at $\leq 0.08$ per task, $\leq 0.001$ per layer, and $\leq 0.49$ per configuration arm. We do not propose a parameter that needs tuning; we show this parameter does not need tuning.

Boundary: concession presupposes sparse critical far tokens. On the synthetic far-dense task (RULER multikey\_3) the far candidates explode and concession fails (far budget granularity dominates; Quest collapses on the same task with the same structure). The length gradient of the gain is clean: the partition versus the single pool is $-0.14$/$-0.25$ at RULER 4K/8K and $+0.91$ at 16K, so partitioning is neutral at short context and amplified at long context: a design constraint for long-context methods. The value of the dynamic signal lies in the layer-skip gate (safety and compute savings), not in quota dispatch: one prefill-side far statistic correlates 0.924 with decode-side far quality, yet the dispatch it drives gains nothing; this is re-verified on the 13 end-to-end configuration arms -- the rank correlation between the offline quality proxy and e2e accuracy is only 0.166 (confidence interval containing 0), and picking arms by the proxy loses 1.18 against the end-to-end optimum, so the proxy signal is actively misleading for quota election, not merely neutral.

\begin{figure}[h]\centering
\includegraphics[width=0.95\linewidth]{figures/fig5_abg_scan.pdf}
\caption{Budget parameter surface ($\alpha/\beta/\gamma$ scan): $\beta$ flat over 0.25--0.75, $\alpha$ monotone to the boundary, oracle upper bounds $+0.08$/$+0.0008$ per task/layer. That the parameter needs no tuning is itself a deployment advantage. Panel (d): offline quality proxy versus e2e accuracy over the 13 end-to-end configuration arms -- the two protocols correlate weakly in rank (a ranking reversal), and configuration election trusts only e2e measurements (see the ablation section).}
\label{fig:abg}
\end{figure}

\subsection{Observation 3: Dimensionality-Reduction Freedom Comes from Position Priors}\label{sec:obsc}

A 32-of-128 (25\%) low-frequency tail subspace loses no indexing quality: selection coverage over the full-dimensional top-1024 is 0.729 versus 0.732 for the same pipeline at $d'{=}128$; a random 32-dimensional subspace collapses to 0.32 and high-frequency dimensions to 0.137. Which 25\% matters more than compressing to 25\%. Inside the tail the two segments complement: the RoPE low-frequency tail alone reaches 0.576 and the nope tail alone 0.387, and their concatenation reaches 0.811 ($+0.24$ combination gain), each contributing what the other lacks. Unsupervised PCA compresses the fine stage further to $d{=}16$ near-losslessly (0.766 versus 0.804), and a kv-head-shared SVD basis matches tail32 on 8/8 samples. The five-arm e2e ablation reproduces the same ranking: rope64 55.98 $>$ full128 54.72 $>$ tail32 54.23 $>$ nope64 52.4 $\gg$ random32 46.01 $\gg$ highfreq32 33.8 (hotpotqa F1).

Three layers of evidence converge on one law: the dimensionality-reduction freedom of a training-free method comes from position priors (distribution-independent structure), not from learning or adaptation. First, arbitrary subspaces collapse while position-structured subspaces are stable: the freedom is not any 32 dimensions but the position-stationary 32. Second, unsupervised spectral structure also works: the far-region content spectrum is concentrated (PCA $d{=}16$ near-lossless) and GQA heads are redundant (a shared projection matches per-head), so a shared projection cache is training-free feasible on pure-GQA models; the MLA latent indexer of DeepSeek-style models is its trained counterpart. Third, the protocol gap settles that adaptation is infeasible: supervised projections hold in trace replay (92--94\%) yet collapse end-to-end, because a projection basis is fragile under decode query-distribution drift while the rotation-pair frequency prior is independent of the query distribution; query-adaptive compression (correlation $\approx 0$, churn 22\%) and nonlinear manifold learning (UMAP/t-SNE, no quality gain at four orders of magnitude higher cost) are both rejected.

Design implication: L1 takes the tail subspace; the full four-quadrant characterization of selection methodology appears in \S\ref{sec:l2}. The dividing line to DSA sits here: DSA trades training (KL-aligning the attention distribution) for selection ability and can learn arbitrary projections; a training-free method can use only distribution-independent position priors. The protocol gap, where trace replay holds but e2e collapses, is the mathematical boundary between these two schools rather than an implementation failure. The universality boundary is reported honestly: on the 32B model the mechanism is present but weaker (low-frequency 32 dimensions at 0.66--0.87, within 0.08 of full dimensionality).

% =====================================================================
% 3 Method: follows the MoBA Section 2 structure: formalize in
% Preliminaries -> architecture + core equations -> two causality designs
% -> design-choice discussion (subspace / budget partition) -> relation to
% sink/SWA -> Implementation.
% =====================================================================
\section{Method}

\subsection{Preliminaries: Standard Attention and the Sparse Selection Problem}

% MoBA 2.1 style: single-head formalization, kept concise.
We first recall standard attention. A single query token $q\in\mathbb{R}^{1\times d}$ attends to $N$ KV tokens ($K,V\in\mathbb{R}^{N\times d}$):
\begin{equation}
\mathrm{Attn}(q, K, V) = \mathrm{Softmax}\!\left(qK^\top\right)V.
\end{equation}
The sparse indexing problem is to select for $q$ a set of positions $I\subseteq[N]$ of size $K_2$ such that $\mathrm{Attn}(q, K[I], V[I])$ approximates the full result. We treat the single-head case for clarity; per-kv-head scoring under GQA and the cross-head unevenness of distant mass are discussed in the partition mechanism of Section~\ref{sec:l2}. Unlike a training-based indexer, a training-free indexer admits no parameters: selection quality depends entirely on which signal scores candidates and how the budget is organized.

\subsection{PSI Architecture}\label{sec:l2}

\begin{figure}[h]\centering
\includegraphics[width=0.85\linewidth]{figures/fig7_tli_architecture.pdf}
\caption{PSI two-level indexer architecture: L1 block-level minmax upper-bound coarse selection (subspace $d'{=}32$ with 4bit quantization) $\rightarrow$ L2 far/near partitioned token-level fine selection $\rightarrow$ a unified sparse attention fused kernel.}
\label{fig:arch}
\end{figure}

% MoBA 2.2 style: Different from standard attention where..., X enables...
Different from standard attention, where every query reads the full context, PSI selects $K_2{=}1024$ KV tokens per query head:
\begin{equation}
\mathrm{PSI}(q, K, V) = \mathrm{Attn}\!\left(q,\ K[I_{\mathrm{L1}\rightarrow\mathrm{L2}}],\ V[I_{\mathrm{L1}\rightarrow\mathrm{L2}}]\right),
\end{equation}
where $I_{\mathrm{L1}\rightarrow\mathrm{L2}}$ is produced by a two-level cascade (Figure~\ref{fig:arch}). The design contribution of PSI is not the cascade form itself (HISA~\cite{hisa} already validated it) but the choices made on two orthogonal dimensions: \textbf{the score representation takes an upper bound}, and \textbf{the budget organization takes a partition}.

\textbf{L1: block-level upper-bound coarse selection.} The context is split into blocks of $B{=}64$. For each block and each kv-head, PSI maintains min/max bounds over a position-stable low-frequency subspace ($d'{=}32$) in 4bit quantized storage (128 dimensions compressed to 40B/token-head). The query score of a block takes, per dimension, the bound endpoint matching the sign of $q_d$ and sums these optimistic terms:
\begin{equation}
s_{\mathrm{blk}}(q, i) \;=\; \sum_{d=1}^{d'} \max\!\big(q_d\, k^{\min}_{i,d},\; q_d\, k^{\max}_{i,d}\big)\;\geq\; \max_{j\in I_i}\; \langle q, k_j \rangle.
\end{equation}
This upper bound guarantees \textbf{no miss within the subspace-score protocol}: the block containing a genuinely high-scoring token scores at least as high as that token on the subspace and always enters the top-$K_1$ (the subspace itself is lossy with respect to the full dot product; that loss is characterized separately by the selection methodology in \S3.3). This contrasts with the page-min lower bound of Quest~\cite{quest}, which underestimates genuinely high-scoring pages and misses more as distractors multiply. The top $K_1{=}128$ blocks proceed to L2.

\textbf{L2: far/near partitioned fine selection.} The sink and the sliding window (the most recent $sw_{lo}$ tokens, whose locality is guaranteed by sliding-window semantics and needs no scoring) form the forced-retention region and consume no partitioned budget; the remaining mid budget splits into two pools. The far pool applies token-level 4bit fine selection over the blocks admitted by L1, taking the far quota $F$ (saturating at 128--256); the near pool fine-selects the near region under the same protocol. The anti-eviction mechanism of the partition is as follows. Under GQA the distant mass is extremely uneven across kv-heads (trace measurement: one head alone holds 0.804 while other heads in the same group are near zero), so in a single-pool top-k the far candidates compete against near high scorers in the same pool and the retrieval-critical tokens of far-heavy layers are systematically evicted. The partition gives far its own quota, and near-side deductions carry a near\_floor guard against boundary regression.

\textbf{Formal definition of the budget.} Let the query prefix length be $S=t+1$, the mid-region length be $L_{\mathrm{mid}}=S-|I_{\mathrm{sink}}|-sw_{lo}$, and the near-region length be $\ell_{\mathrm{near}}=\lceil\alpha L_{\mathrm{mid}}\rceil$. The two pools and the final selection set are
\begin{align}
I_{\mathrm{far}} &= \operatorname*{TopK}_{F}\Big\{\langle q,k_j\rangle:\ j\in\bigcup\nolimits_{i\in I_{\mathrm{L1}}} I_i\cap\big[\,|I_{\mathrm{sink}}|,\ S-sw_{lo}-\ell_{\mathrm{near}}\,\big]\Big\},\\
I_{\mathrm{near}} &= \operatorname*{TopK}_{K_2^{\mathrm{mid}}-F}\Big\{\langle q,k_j\rangle:\ j\in\big[\,S-sw_{lo}-\ell_{\mathrm{near}},\ S-sw_{lo}\,\big]\Big\},\\
I_{\mathrm{L1}\rightarrow\mathrm{L2}} &= I_{\mathrm{sink}}\ \cup\ \big[\,S-sw_{lo},\ t\,\big]\ \cup\ I_{\mathrm{near}}\ \cup\ I_{\mathrm{far}},
\end{align}
where the mid budget is $K_2^{\mathrm{mid}}=K_2-|I_{\mathrm{sink}}|-sw_{lo}$ and the far quota is $F=\max\big(64,\ K_2^{\mathrm{mid}}-\lfloor\gamma\,\lceil\ell_{\mathrm{near}}/B\rceil B\rfloor\big)$: $\gamma$ converts the near block count into the near fine-selection quota, far takes the remaining budget with a floor of 64, and $\beta$ acts on the near page budget at L1.

\textbf{Two designs that preserve causality.} An autoregressive language model must not let a query access future tokens. PSI preserves causality through two concrete designs. First, \textbf{pool boundaries are causal boundaries}: the far pool's upper position bound is $\leq t+1-\ell_{\mathrm{near}}-sw_{lo}$ and the near pool stays below $sw_{lo}$, so the union of the two pools falls inside the causal range by construction, with no explicit causal mask (garbage blocks from L1 also fall outside the pools). Second, \textbf{uniform coverage for early rows}: within the first prefill chunk, early rows whose position $t_r$ is below $K_2$ lack candidates, and if sentinel or future positions returned by topk leak in, non-causal leakage propagates layer by layer (we observed repetitive-loop generation collapse on a 30B model with exactly this root cause; the 8B model tolerated the same bug without collapsing). Early rows use a uniform repeated grid, which is equivalent to dense.

\textbf{Relation to sink and sliding-window attention.} Sink tokens (StreamingLLM~\cite{streamingllm}) and sliding-window attention (Longformer, BigBird~\cite{longformer,bigbird}) are two established static sparse structures, and we adopt them as known components: sink tokens carry 0.37--0.71 of attention mass and the sliding window carries near locality, together forming the forced-retention region of PSI that does not consume the partitioned budget. Our contribution is not the retention region itself but \textbf{the allocation of the remaining budget between far retrieval and near fine selection}. As corroborating evidence, we measured in baseline adaptation that SnapKV, H2O, and PyramidKV without forced sink retention emit EOS as their first token on Qwen3-8B (the mean sink weight in the last layer is 0.592 while shallow-layer selection signals do not point to sink, a layer-by-layer progressive distortion): any sparse method that does not handle sink explicitly distorts on this model family.

\subsection{Design Choice Discussion}

% MoBA style: Next, we discuss some additional key design choices...
We next discuss three key design choices of PSI.

\textbf{Subspace choice.} L1 takes the 32-dimensional low-frequency tail because the indexing signal must preserve the ranking power of full-dimensional scores after reduction. The rotate\_half RoPE of Qwen3 makes dot products over the low-frequency dimensions rank-consistent with full-dimensional dot products (selection coverage $C{=}0.729$ over the full-dimensional top-1024, versus 0.732 for the same pipeline at $d'{=}128$ (4bit quantization and cascade losses included)), while a random subspace collapses to 0.32 and high-frequency dimensions to 0.137. The mechanism is confirmed by a triple verdict (Figure~\ref{fig:dim}): all 128 dimensions of Qwen3 participate in RoPE, and under the rotate\_half layout frequency $j$ acts on the dimension pair $(j, j{+}64)$, so tail32 (dimensions 48..63 and 112..127) equals \textbf{the 16 complete lowest-frequency rotation pairs}. The dot product decomposes by frequency as $\sum_j [(q_j k_j + q_{j+64} k_{j+64})\cos(\Delta\,\omega_j) + (q_j k_{j+64} - q_{j+64} k_j)\sin(\Delta\,\omega_j)]$: at large distant $\Delta$ the low-frequency terms are smooth and highly correlated while the high-frequency terms oscillate and cancel, so the low-frequency pairs carry the long-range retrieval signal. Three verdicts support this structural account: frequency monotonicity ($R_{\mathrm{far}}$: low-frequency 16 pairs 0.811, mid-frequency 0.360, high-frequency 0.004); pairing completeness (a misaligned pairing, taking the low-frequency 16 from the first half but a mid-frequency source for the second half, collapses to 0.567 with nearly identical segment positions); and saturation at 16 pairs (16 pairs reach 0.811, near-saturated, while 8 pairs reach only 0.473). The direction replicates on the 32B model (complete pairs 0.829 $>$ misaligned 0.645 $>$ single half 0.54), with the saturation width shifting to 24 pairs (0.904).

The full characterization of selection methodology (Table~\ref{tab:select}): crossing static (index-compatible) with adaptive (per-query) selection yields four quadrants. Per-dimension spectral top-32 reaches only 0.759: the two elements of a rotation pair are entangled in the dot product, so per-dimension importance cannot in principle recover the pairing structure. SparQ-style~\cite{sparq} per-query $|q|$-magnitude top-32 reaches 0.848, and pair-aware adaptive selection (ranking complete pairs by $|q_j|+|q_{j+64}|$) reaches 0.865. Adaptive selection does yield a small gain, but its subspace changes per query, which is structurally incompatible with an incremental minmax block index built before the query arrives (SparQ pays a per-query dimension-gather cost for exactly this reason and builds no index). The intermediate point is per-layer static selection: the top-16 complete pairs by pair magnitude from the $q$ statistics of the last prefill segment, fixed per layer and indexable, at 0.849, which keeps seven tenths of the adaptive gain at 3.8 points above tail32. \textbf{tail32 is precisely the optimum of the zero-signal-cost feasible region}, not the global optimum. Static pair selection is structurally indexable, and we have run its e2e verification: the static-pair arm (same configuration, selecting pairs from prefill $q$ statistics) gains $+3.8$ points in replay yet loses $-0.50$ at e2e (50.04 versus 50.54, musique $-5.27$) -- a selection basis derived from prefill statistics is structurally homologous to a trained projection and fails under decode query-distribution drift, so replay gains do not extrapolate. Trained projections rank more accurately (92--94\% in trace replay) yet collapse end-to-end: the projection basis is fragile under decode query-distribution drift, whereas the rotation-pair frequency prior is independent of the query distribution and stays stable. This protocol gap, where the trace protocol holds and the e2e protocol collapses, is an empirical characterization of the dividing line between the training-free and the training-based families.

\begin{figure}[h]\centering
\includegraphics[width=0.95\linewidth]{figures/fig2_dim_reduction_story.pdf}
\caption{Dimensionality-reduction freedom comes from position priors: (a) tail32 equals the 16 complete lowest-frequency rotation pairs (frequency monotonicity, misaligned pairing collapses to 0.567, saturation at 16 pairs); (b) a shared SVD basis at $d{=}16$ matches tail32 on 8/8 samples; (c) linear beats nonlinear, shared beats per-head -- UMAP was evaluated on only 2 samples (fit cost 80--167s per layer-sample), and the paired reference bar gives tail32 on the same 2 samples (0.835 versus 0.813).}
\label{fig:dim}
\end{figure}

\begin{table}[h]\centering
\caption{Four quadrants of subspace selection methodology (far-region mass recall $R_{\mathrm{far}}$, mean over 8 samples).}
\label{tab:select}
\begin{tabular}{llcc}
\toprule
Selection method & Signal & recall & Index-compatible \\
\midrule
Per-dim spectral top-32 & data-driven (static) & 0.759 & compatible \\
tail32 & frequency prior (static) & 0.811 & compatible \\
SparQ-style $|q|$ top-32 & query-adaptive & 0.848 & incompatible \\
Pair-magnitude top-16 pairs & query-adaptive & 0.865 & incompatible \\
Per-layer static pairs & prefill $q$ statistics & 0.849 & compatible \\
\bottomrule
\end{tabular}
\end{table}

\textbf{Budget partition parameters.} The budget space is described by three parameters: $\alpha{=}0.125$ (near-region fraction), $\beta{=}0.375$ (near page budget), and $\gamma{=}0.625$ (near token budget discount), with far/near methods set to (minmax, avg). The parameter surface is characterized in full by a dual-metric grid (methodology in the ablation section, \S\ref{sec:abl}): the offline quality grid scans a $9\times9\times9$ $(\alpha,\beta,\gamma)$ cube (0.125 step, 3600 valid configurations after legality-constraint filtering, 16-sample trace means) over five method combinations, and the e2e grid curates 13 configuration arms -- deduplicated by effective e2e budget from each combination's quality-optimal configurations -- measured on the two LongBench screen tasks (the full list in Table~\ref{tab:e2egrid}). Three regularities: $\alpha$ is monotone toward its boundary (the less far budget the better, limited by near\_floor); $\beta$ is flat over 0.25--0.75 (multi-hop tasks favor 0.375); and $\gamma$ must scale proportionally with $K_2$ (the lesson of $\gamma{=}0.5$ under $K_2{=}1024$, where far retains only the floor quota of 64 and F1 collapses to 13.51; note also that the 2$\times$ gap between the offline total budget $B_{\mathrm{TOK}}{=}2048$ and the e2e $K_2{=}1024$ lets $\gamma$ be truncated and flattened on the e2e side by the near fine-selection quota, so cross-protocol arm selection must deduplicate by effective configuration). The most important finding is the flatness itself: the per-task oracle that perfectly selects $\beta$ gains only +0.08, the per-layer oracle only +0.0008, and the per-configuration-arm oracle $+0.21$ to $+0.49$, so \textbf{the budget partition is insensitive to configuration error and needs no tuning}. The main-table configuration is elected by the e2e grid (the combination ranking of the offline quality protocol reverses at e2e, and the election trusts only e2e measurements; see \S\ref{sec:abl}). This is the deployment-simplicity argument: against the 1000-step warm-up of DSA, one fixed PSI configuration covers 12/13 tasks (the LongBench real-task profile; the synthetic-needle profile takes the reference arm as a two-way selection, see \S4.2 and the limitations).

\textbf{Gate threshold.} The dynamic gate uses last-segment prefill far statistics to identify layers whose far mass is near zero. The verdict on the threshold $\tau$ comes from a two-task gradient: at $\tau{=}0.005$, musique moves +0.23 and qasper $-0.41$; at $\tau{=}0.015$, +0.27 and $-0.67$. The gradient is monotone, and the real triggering regime takes small losses without collapsing (against 4.8--5.9 lost by the static variant). The gate is therefore positioned as a conservative switch that saves compute safely.

\subsection{Implementation}

% MoBA 2.3 style: We provide a high-performance implementation..., consisting of N major steps.
We implement PSI in the production form of sglang; the selection chain consists of four kernel components:
\begin{itemize}
\item \textbf{Fused L1 block scoring}: one Triton kernel performs gather, 4bit dequantization (shared scale table), GEMV scoring, and block-level reduction, 3.6$\times$ on the 32K shape;
\item \textbf{L2 cascaded dual kernel}: a single launch performs far/near dual-pool scoring and partitioned topk, with pool boundaries doubling as causal boundaries (replacing the eager fine-matrix causal mask);
\item \textbf{Unified sparse attention fused kernel}: gather, online softmax, and bf16 writeback, 11--16$\times$ at the kernel level;
\item \textbf{O(n) exact incremental indexing}: per-token incremental block-bound updates match full rebuild bit for bit, at 0.128ms independent of sequence length.
\end{itemize}
Three system-side pieces connect the design to production: paged addressing (selection outputs logical positions resolved through the req\_to\_token indirection into the KV pool, zero-copy); a per-layer shared, pre-allocated index pool (the launch count of batched decode selection is independent of batch size); and a CUDA-graph capture region with zero host synchronization. The design constraints come from the H20 profile: Tensor Core throughput is only 15\% of H100 while HBM bandwidth is in the same class, so the first-order constraint on scoring kernels is bandwidth rather than compute. We measured and rejected both the Tensor Core path and the TMA transposed layout (arithmetic-intensity profiling shows the pure gather path is already bandwidth-saturated), so the kernel line of attack is eliminating intermediate materialization and merging launches.

% =====================================================================
% 3 Experiments: follows the MoBA Section 3 pattern (goal sentence ->
% experiment -> numbers -> mechanism -> "This suggests that...").
% =====================================================================
\section{Experiments}

\subsection{Metrics}

\textbf{End-to-end quality.} RULER uses the official string\_match\_all $\times$100; LongBench averages each task's official metric (F1 / Rouge / accuracy).

\textbf{Attention quality (mass-type metrics).} For decode step $t$ and layer $\ell$, let $a_j$ denote the full-dimensional dense attention weight, the softmax over $\langle q, k_j \rangle / \sqrt{d}$. Two protocols are used:
\begin{itemize}
\item \emph{Selection coverage} ($C$): with the top-$K$ reference token set $M^\ast$ selected by full-dimensional scoring as the denominator, $C(M) = \sum_{j \in M} a_j / \sum_{j \in M^\ast} a_j$ measures how much information the selected set retains when subspace scoring replaces full-dimensional scoring. The 0.729 in Contribution \#1 uses this protocol ($K{=}1024$).
\item \emph{Far-region mass recall} ($R_{\mathrm{far}}$): the denominator is restricted to the far region (tokens outside sink and sliding window), $R_{\mathrm{far}}(M) = \sum_{j \in M \cap F} a_j / \sum_{j \in F} a_j$, measuring how much of the distant retrieval signal the sparse selection captures. Both the 0.811 in the triple verdict and the 0.804 in Ablation A use this protocol; the numerical difference comes from different experiment batches (the former is the 8-sample mean of the rotation-pair verdict experiments, the latter the 8-sample mean of the dimensionality-reduction experiments, with partially non-overlapping sample sets); each is internally compared against the same-batch tail32 baseline and is not compared across batches.
\end{itemize}
All mass metrics are aggregated with GQA kv-head-level distribution weighting (the $a_j$ of query heads within a group are computed separately before aggregation), preventing head averaging from masking the head concentration of far mass.

\subsection{End-to-End Quality}

% MoBA 3.1 style: In this section, we conduct... to validate...
This section validates the quality of PSI on two benchmarks: RULER~\cite{ruler} (3 lengths $\times$ 11 tasks $\times$ n=100, NVIDIA's official pre-generated set) probes long-context retrieval, and 13 LongBench~\cite{longbench} subsets probe the mixed regime of real tasks. Among the compared methods, TIA is the first-generation training-free block upper-bound indexer built on the same pipeline as ours (block min/max upper bounds with 4bit subspace quantization, no partitioned budget and no perception gate) and serves as the most direct structural ablation reference for PSI; Quest represents the lower-bound family; the three KV-compression baselines run in the unified monkeypatch pipeline. All experiments use the real weights of Qwen3-8B~\cite{qwen3}, and PSI uses its final configuration (the strict protocol in which sink and sliding-window forced-retention tokens consume no budget). The L1 block upper bound of the main-table arm is computed on the full 128 dimensions (without 4bit quantization), while the L2 token-level fine selection runs on the 4bit-quantized tail32 subspace. The main-arm configuration was additionally re-measured with a tail32 L1 bound across all 13 tasks (50.53 versus 50.78, paired difference $-0.25$, 95\% CI $[-0.57, +0.06]$ containing 0; per-sample bootstrap, $B{=}10000$): the L1 bound dimension is statistically indistinguishable on the main arm -- yet a directional structure exists, with far-retrieval-heavy hotpotqa dropping 1.81 points ($\gamma{=}0.625$'s tight far budget amplifies the bound's information loss) while code and summarization tasks stay within 0.3 points; the earlier paired calibration on the $\gamma{=}0.125$ reference arm (54.83 versus 54.43, 33.22 versus 34.76) is consistent with this: dimension sensitivity grows as the far budget tightens. The main table keeps the full-dimensional L1 protocol, with tail32-L1 as a same-budget ablation arm (its difference against FullKV is $+0.16$ under the per-sample bootstrap-mean protocol, with a CI containing 0, so the parity claim holds under both L1 protocols).

\begin{table}[h]\centering
\caption{RULER final table (string\_match\_all $\times$100; PSI uses the partitioned configuration: (minmax, avg) scoring with $\alpha{=}0.125$, $\beta{=}0.25$, $\gamma{=}0.125$, $K_2{=}1024$, the reference configuration optimal for the synthetic-needle profile -- the LongBench main-table arm uses $\beta{=}0.375/\gamma{=}0.625$, elected by the end-to-end grid (\S\ref{sec:abl}), and its same-table RULER rescan lands at an overall 85.83 (the $-2.10$ gap concentrates on the cwe/fwe word-list tasks and extrapolates to $-2.15$ at 32K, stable with length; the remaining tasks near-parity, 7 exact and multikey\_3/multivalue within 3 points; see the text); PSI-single = the same-budget $\alpha{=}0$ single-pool ablation arm; the 32K$^\dagger$ row uses data produced with NVIDIA's official generator (essay corpus same-source with 4K--16K, n=100, with the niah\_single tasks at full marks for the upper-bound family as a pipeline health check, Quest at 91--98), where most tasks enter a family-wide degradation regime (see the text); in that row TIA and PSI-single score identically on all 11 tasks -- the two command paths under the same pipeline and budget produce byte-identical outputs (verified by independent dual-GPU runs), an algorithmic-equivalence reproduction rather than a transcription error; the Overall AVG keeps the 4K--16K three-length protocol).}
\begin{tabular}{llllll}
\toprule
 & FullKV & TIA@1024 & PSI-single & PSI@1024 & Quest@1024 \\
\midrule
4K & 91.59 & \textbf{91.63} & 91.63 & 91.49 & 87.91 \\
8K & \textbf{89.14} & 88.90 & 88.87 & 88.62 & 81.76 \\
16K & \textbf{85.41} & 84.53 & 82.77 & 83.68 & 70.18 \\
32K$^\dagger$ & 59.38 & 60.77 & 60.77 & 60.78 & 47.63 \\
\addlinespace
Overall AVG & \textbf{88.71} & 88.35 & 87.76 & 87.93 & 79.95 \\
\bottomrule
\end{tabular}
\label{tab:ruler}
\end{table}

\begin{table}[h]\centering
\small\setlength{\tabcolsep}{4pt}
\caption{RULER 16K per-task breakdown (string\_match\_all $\times$100; configurations as in the table above). single\_1--3 and multikey\_1 are 100 for all five methods and multiquery is 25 for all including FullKV (model ceiling), so they are omitted.}
\label{tab:ruler16k}
\begin{tabular}{lccccc}
\toprule
 & FullKV & Quest & TIA & PSI-single & PSI@1024 \\
\midrule
multikey\_2 & 100.0 & 60.0 & \textbf{100.0} & \textbf{100.0} & \textbf{100.0} \\
multikey\_3 & \textbf{100.0} & 21.0 & 89.0 & 75.0 & 80.0 \\
multivalue & \textbf{44.25} & 41.25 & \textbf{44.25} & \textbf{44.25} & \textbf{44.25} \\
cwe & \textbf{83.60} & 29.50 & 78.90 & 73.20 & 78.60 \\
fwe & 86.67 & \textbf{96.67} & 92.67 & 93.00 & 92.67 \\
vt & \textbf{100.0} & 98.60 & \textbf{100.0} & \textbf{100.0} & \textbf{100.0} \\
16K AVG & \textbf{85.41} & 70.18 & 84.53 & 82.77 & 83.68 \\
\bottomrule
\end{tabular}
\end{table}

Three observations hold as shown in Table~\ref{tab:ruler16k}. First, the lead of PSI over Quest grows with length (+3.58 $\rightarrow$ +6.86 $\rightarrow$ +13.50) and holds at +13.15 in the 32K family-wide degradation regime: at 16K, Quest collapses on the multi-distractor retrieval tasks (multikey\_3 21, multikey\_2 60, cwe 29.5) while PSI stays within the correct order of magnitude (multikey\_2 and vt at 100, cwe 78.6), because lower-bound misses worsen as distractors multiply whereas upper-bound coarse selection has no such degradation. Second, the gain of the partition over the single-pool ablation arm turns positive only at 16K ($-0.14 \rightarrow -0.25 \rightarrow +0.91$, all measured against the same-configuration single-pool arm): ultra-long context is the payoff regime of partitioned far retrieval, and the 16K partition gain concentrates in multikey\_3 (75.0 $\rightarrow$ 80.0, +5.0) and cwe (73.2 $\rightarrow$ 78.6, +5.4). Third, absolute scores must be read against FullKV at the same length: aggregate enumeration tasks degrade for FullKV itself as length grows (multivalue 89.5 $\rightarrow$ 69.25 $\rightarrow$ 44.25 is the enumeration ceiling of the model), so the same-length gap is the sparse effect. The same-table rescan of the main arm $\gamma{=}0.625$ (a full 33-arm sweep) delivers the dual-arm verdict: an overall 85.83 (89.77/86.28/81.45 at 4K/8K/16K), $-2.10$ below the reference arm, with the damage concentrated in the word-list aggregation task cwe (82.0/68.6/55.5, its absolute score degrading over 4K--16K) and fwe ($-3.0$/$-1.67$/0.0) while 7 of the other 9 tasks (the three single tasks, multikey\_1/2, multiquery, and vt) match the reference arm exactly, with multikey\_3 ($-3.0$/$-1.0$) and multivalue ($-0.5$/$-1.0$/$-0.5$) taking small losses: the tight far budget of $\gamma{=}0.625$ fails on synthetic word-list tasks with exploding far candidates, isomorphic to the far-dense boundary criterion (multikey\_3 $-2.00$). Both arms are reported as the deployment protocol: the real-task profile takes the main arm (LongBench 50.78), the synthetic-needle profile takes the reference arm (87.93), consistent with the two-profiles-each-take-their-arm reading of the $\beta$ cross table (Table~\ref{tab:betacross}). The 32K supplementary point (six arms on the same table, Table~\ref{tab:ruler}) extrapolates the dual-arm gap to $-2.15$ -- matching $-2.10$ exactly, so the $\gamma$ damage is stable with length rather than amplifying, with cwe at $-18.2$ accounting for 77\%; most tasks of that length enter a family-wide degradation regime (multikey\_3 at 4--16 for the whole family, fwe with FullKV at 7.0 below the sparse arms, vt at 31--35 for the whole family), absolute scores are read against same-length FullKV, and the only true sparse differentiation point that keeps its quality is again cwe (FullKV/TIA/single-pool at 83.5, the PSI reference arm at 80.1, the main arm at 61.9, Quest at 24.5) -- the upper-bound-family-versus-Quest contrast is at its cleanest at 32K.

\begin{table}[h]\centering
\small
\caption{LongBench 13-subset final table (official per-task metric averaged; PSI = the partitioned configuration (minmax, avg) with $\alpha{=}0.125$, $\beta{=}0.375$, $\gamma{=}0.625$, elected by the end-to-end leg of the dual-metric grid (\S\ref{sec:abl}); PSI-single-pool = the same-budget $\alpha{=}0$ single-pool ablation arm; the three KV-compression baselines run in the unified monkeypatch pipeline at a 1024 budget, without sink-guard adaptation; StreamingLLM = budget 1024 in the unified pipeline; MoBA = unified-harness reproduction (full-dimensional chunk-mean gate selecting the top-16 blocks fully expanded, sink/swa retention and budget matched to the same protocol).}
\label{tab:longbench}
\small\setlength{\tabcolsep}{1.5pt}
\begin{tabular}{lcccccccccc}
\toprule
 & FullKV & Quest & TIA & PSI-single & PSI & SnapKV & H2O & PyramidKV & StreamLLM & MoBA \\

\midrule
hotpotqa & 53.48 & 45.74 & 53.89 & 54.93 & 55.44 & 23.64 & 2.68 & 25.10 & 2.90 & 51.02 \\
2wikimqa & 38.29 & 38.46 & 38.27 & 38.57 & 39.32 & 28.62 & 7.90 & 26.01 & 5.88 & 37.16 \\
musique & 32.14 & 27.25 & 32.28 & 29.74 & 34.71 & 8.61 & 0.63 & 8.99 & 1.99 & 29.74 \\
passage\_retrieval\_en & 100.00 & 98.50 & 99.50 & 99.50 & 99.00 & 4.00 & 2.50 & 4.50 & 0.25 & 94.50 \\
qasper & 44.17 & 40.13 & 44.03 & 44.31 & 43.95 & 16.84 & 10.18 & 18.58 & 8.69 & 42.44 \\
multifieldqa\_en & 53.40 & 51.01 & 53.19 & 52.87 & 53.05 & 20.24 & 13.45 & 23.68 & 9.95 & 52.83 \\
gov\_report & 33.17 & 32.13 & 33.43 & 33.49 & 33.52 & 19.21 & 17.23 & 20.56 & 28.85 & 32.79 \\
qmsum & 23.53 & 22.21 & 23.90 & 23.87 & 23.58 & 18.46 & 8.82 & 18.54 & 15.84 & 23.12 \\
multi\_news & 24.93 & 24.95 & 24.73 & 24.64 & 24.44 & 21.04 & 22.79 & 21.37 & 24.56 & 24.62 \\
narrativeqa & 25.61 & 20.64 & 22.18 & 24.46 & 26.31 & 11.53 & 1.41 & 12.03 & 2.95 & 23.70 \\
triviaqa & 90.71 & 87.55 & 89.82 & 89.88 & 90.57 & 69.77 & 18.01 & 75.03 & 14.30 & 90.46 \\
lcc & 68.81 & 68.34 & 69.14 & 69.28 & 69.28 & 67.19 & 54.06 & 67.05 & 36.94 & 68.98 \\
repobench & 66.50 & 63.41 & 66.40 & 66.58 & 66.97 & 50.93 & 33.65 & 53.17 & 31.44 & 68.13 \\
Overall & 50.36 & 47.72 & 50.06 & 50.16 & \textbf{50.78} & 27.70 & 14.87 & 28.82 & 14.20 & 49.19 \\
\bottomrule
\end{tabular}
\end{table}

On the 13 LongBench subsets (Table~\ref{tab:longbench}), PSI scores \textbf{50.78}, exceeding FullKV (50.36) by $+0.42$, Quest by $+3.06$, the first-generation TIA of the same pipeline by $+0.72$, and the training-based gate MoBA by $+1.59$. The significance verdict is reported as measured: the 95\% confidence interval of $+0.42$ is $[-0.18, +1.02]$ and contains 0 (per-sample paired bootstrap, $B{=}10000$, resampling within tasks and aggregating with equal weights, with PSI and FullKV matched sample-by-sample), and the sign test gives 8 wins and 5 losses at $p{=}0.58$ -- the accuracy claim against FullKV is parity-level rather than a statistically significant lead, and the paper is accordingly positioned on the speed win at matched accuracy. The gain concentrates on far-heavy multi-hop tasks: on musique, PSI reaches 34.71, the highest among all methods including FullKV (+2.57 over FullKV, +4.97 over the same-budget single-pool arm; the single-task 95\% CI $[-1.29, +6.37]$ also contains 0), which is a structural signal of partitioned far retrieval; across the 13 tasks PSI wins 8 and loses 5 against FullKV ($p{=}0.58$), the positive side driven by musique $+2.57$, hotpotqa $+1.96$, 2wikimqa $+1.03$, and narrativeqa $+0.70$, with all losses at most 1 point (passage\_retrieval at $-1.00$ is the largest loss, followed by multi\_news $-0.49$ and multifieldqa $-0.35$). The main-table arm is elected by the dual-metric grid of \S\ref{sec:abl} ($\gamma{=}0.625$) and gains $+0.24$ over the old deployment-reference arm ($\gamma{=}0.125$, 50.54), with the gain coming from narrativeqa $+2.48$ and hotpotqa $+1.01$, while musique and qasper dip slightly ($-0.05$/$-0.56$); the election is run on hotpotqa and musique, which are in-sample, and the held-out 11 tasks differ from FullKV by $+0.08$ (95\% CI $[-0.46, +0.60]$) -- the bulk of the $+0.42$ difference comes from the two in-sample election tasks, and the full-set claim should be read as parity. The full collapse of the three KV-compression baselines confirms the sink claim: SnapKV lands at 27.70 and H2O collapses further (musique 0.63, hotpotqa 2.68), because without an explicit sink guard Qwen3-8B emits EOS as its first token (mechanism at the end of \S3.2); only the code tasks survive partially (lcc 67.19 and 54.06). The static-sparse control completes the same conclusion: StreamingLLM (sink plus sliding window, budget 1024) lands at an overall 14.20 (Table~\ref{tab:longbench}), with systematic collapse on retrieval tasks (passage\_retrieval 0.25, musique 1.99) -- when the budget outside the forced region is zero, far retrieval has nothing to draw on, which delimits from the negative side the working surface that the PSI partition budget must cover. As the representative of training-based gating, MoBA runs in the unified harness (full-dimensional chunk-mean gate selecting the top-16 blocks fully expanded, sink/swa retention and budget matched to the same protocol) at a 13-task full average of 49.19 (Table~\ref{tab:longbench}), $-1.17$ below FullKV; the structural split is clean: QA and retrieval tasks trail the main-table PSI arm by more than 4 points (hotpotqa 51.02 versus 55.44, musique 29.74 versus 34.71), while the gap narrows on code tasks (MoBA leads repobench 68.13 versus 66.97 by 1.16; PSI keeps lcc ahead at 69.28) -- the unbounded chunk gate drops a tier in token-level quality, and its aggregate loss is partly offset by the block-locality bonus on code tasks. The task-profile dependence of the budget mix is given in full as a $\beta \in \{0.25, 0.375\} \times$ \{LongBench, RULER\} cross table (Table~\ref{tab:betacross}, the $\gamma{=}0.125$ reference-configuration family): LongBench multi-hop tasks favor $\beta{=}0.375$ (50.54 versus 50.41; the e2e-grid protocol agrees in direction, 45.08 versus 44.86), and RULER synthetic needle tasks favor $\beta{=}0.25$ (87.93 versus 87.59) -- the diagonal wins, yet all four cells sit within 0.35 of each other, consistent with the parameter-surface flatness in the ablation; the two task profiles each take their arm, and practitioners select by task profile. The per-task-family decomposition gives the boundary criterion for the $\beta$ concession: it holds across all real-task families (musique, the most far-heavy task, gains the most at +1.94) and fails on the synthetic far-dense family (multikey\_3 $-2.00$, where the exploding far-candidate count makes budget granularity dominate): far-dense retrieval is the hard region for the entire sparse-indexer family (Quest collapses on the same task with the same structure), not a defect specific to the $\beta$ mix.

\begin{table}[h]\centering
\caption{$\beta$ cross table over both benchmarks (all other partition parameters held: $\alpha{=}0.125$, $\gamma{=}0.125$, $K_2{=}1024$; full LongBench-13 and RULER-33; both arms belong to the $\gamma{=}0.125$ deployment-reference family -- the election of the $\gamma{=}0.625$ main-table arm appears in \S\ref{sec:abl}).}
\label{tab:betacross}
\begin{tabular}{lcc}
\toprule
 & $\beta{=}0.25$ & $\beta{=}0.375$ \\
\midrule
LongBench-13 & 50.41 & \textbf{50.54} \\
RULER-33 & \textbf{87.93} & 87.59 \\
\bottomrule
\end{tabular}
\end{table}

\subsection{Ablation}\label{sec:abl}

% MoBA style: goal sentence -> numbers -> mechanism -> conclusion.
\textbf{Subspace (A).} The low-frequency $d'{=}32$ reaches a selection coverage of $C{=}0.729$ against the full-dimensional top-1024, versus 0.732 for the same pipeline at $d'{=}128$ (4bit quantization and cascade losses included) (random 0.32, high-frequency 0.137), so the position-stable subspace exists. The mechanism is confirmed by the triple verdict: tail32 equals the 16 complete lowest-frequency rotation pairs (frequency monotonicity $R_{\mathrm{far}}$ 0.811/0.360/0.004, misaligned pairing collapses to 0.567, saturation at 16 pairs). The four-quadrant characterization of selection methods (Table~\ref{tab:select}) bounds the feasible region: per-dimension spectral selection at 0.759 and trained projections both lose to pairing entanglement, query-adaptive pair selection at 0.865 is the oracle ceiling but cannot be indexed, per-layer static pairs at 0.849 are index-compatible yet reverse at e2e ($+3.8$ points in replay, $-0.50$ at e2e, musique $-5.27$, the same protocol gap as trained projections), and tail32 is the e2e-verified zero-signal-cost optimum. A five-arm e2e ablation (same configuration, only the $d'$ selection changed) reproduces the same ranking: rope64 (the first 64 rotation dimensions) at hotpotqa 55.98 / musique 32.96 $>$ full128 54.72/32.82 $>$ tail32 54.23/32.82 $>$ nope64 52.40/30.80 $\gg$ random32 46.01/24.08 $\gg$ highfreq32 33.80/14.26. The pure dimensionality gain of full128 over tail32, $+0.49$ on hotpotqa, shows that 4bit-quantized tail32 already approaches full-dimensional information -- a 4$\times$ dimensionality cost buys under 1 point, so tail32 is the cost-quality balance point and rope64 the quality ceiling reference. The dimensionality-reduction verdict has three parts: unsupervised PCA at $d{=}16$ (a shared SVD basis across kv-heads) is near-lossless ($R_{\mathrm{far}}$ 0.766 versus same-batch tail32 0.804) and is the only recommended option (a per-head basis drops to 0.752, see Figure~\ref{fig:dim}(c)); trained projections hold in trace replay and collapse in e2e (the protocol gap); and nonlinear methods (UMAP, t-SNE) add no quality while costing 80--167s per layer-sample to fit versus 0.128ms per token for incremental indexing, four orders of magnitude and therefore infeasible. A shared SVD basis across kv-heads at $d{=}16$ matches tail32 on 8/8 samples and beats per-head bases, so GQA head redundancy makes the training-free proxy of an MLA-style single-projection cache viable.

\textbf{Partition.} Token-level fine selection in the far region approaches the oracle (0.999 on far-heavy layers); cluster representatives show no consistent advantage under a strict token budget (block scatter-amax is the worst at 0.09--0.39), and cluster centers are not upper bounds, which leaves a theoretical miss hole, so clustering is demoted to a negative result. Partition versus single pool holds on both e2e benchmarks (LongBench +0.62 = 50.78$-$50.16, verifiable from Table~\ref{tab:longbench}; RULER +0.17 = 87.93$-$87.76, the $\beta{=}0.25$ matched arm; both are internal ablation-arm comparisons with large task-level variance, read as structural signals rather than statistical significance), and its disagreement with the trace-replay protocol (the single pool wins slightly at wide budgets) pins down the mechanism: \textbf{the partition gain is budget allocation under a tight budget, not candidate quality}. To rule out the explanation that the trace-replay budget was not tight enough, we rescanned mono versus partition across the full range $B_{\mathrm{TOK}}\in\{256,512,768,1024,2048\}$ under tight budgets: in the trace-mass protocol mono wins slightly at every level ($\Delta\approx-0.003$, partition win rate 0--1/16), so the tight-budget hypothesis is rejected in the trace protocol and the disagreement is metric-level rather than budget-level: mass coverage is masked by the sink/swa/near saturation region and cannot see far-retrieval quality, so the partition gain can only develop on the e2e task structure (musique $+2.43$ overtaking TIA, the RULER 16K $+0.91$ length gradient). The top-$\sigma$ hybrid arm adds third-hand corroboration of partition necessity: abandoning two-level fine selection on the far side for threshold selection degrades e2e monotonically with $\sigma$ (54.31 $\rightarrow$ 52.48), while the same change on the near side is elastic (hotpotqa 55.12, the highest of all configurations, full-dimensional L1 bound protocol) -- far retrieval quality must be protected by block upper-bound coarse selection plus token-level fine selection, while near can relax on the sliding-window floor: the mechanistic basis of the asymmetric partition budget. An ablation of kv-head-level sharing versus per-q-head independent selection (E103, main-arm configuration, both screen tasks) delivers a two-sided verdict: the per-q-head arm at the same $K_2{=}1024$ (each q-head independently selecting 1024, index size $\times$4 at equal compute) scores 53.79/34.51 on hotpotqa/musique against 55.44/34.71 for sharing, a loss of $-1.65$/$-0.20$ -- the within-group mean aggregation is a denoiser, not an information bottleneck; the same-total-index-entry arm (32$\times$256 = 8$\times$1024) collapses to 28.37/12.00, because per-heading the forced region multiplies the sink/swa overhead by $G$ and consumes the entire mid quota ($K_2{=}256$ leaves mid at zero, while the shared arm's forced region takes only 8$\times$256 and leaves 6144 mid entries) -- sharing spreads the forced region across the kv-head level, which is both the source of the index-size saving and a structural gain in budget efficiency.

\textbf{Layer skipping (D').} The offline mean profile skips 13/36 layers (precision 0.92--1.00), but a global static mask does not transfer across tasks (multi-hop tasks lose 4.8--5.9); the dynamic prefill signal correlates at 0.924 and blocks the reverse error, and the real triggering regime takes small losses without collapsing, so the gate is a safety switch. The benefit side has a conservative quantification: the $\tau{=}0.05$ perception arm skips 14.3\% of layers (24/168) at a far-mass loss of 0.8\%, saving selection-chain kernel compute in equal proportion (about 4.2 ms per token step out of a 29.5 ms full chain at S=8K); the aggressive $\tau{=}0.1$ arm skips 30.9\% but the mass loss rises to 2.3\% and touches far retrieval. The upper bound is the offline profile's 13/36 = 36.1\% of layers.

\textbf{The rejection of adaptive $\alpha/\beta$.} Per-request online dispatch of $\alpha/\beta$ is rejected (No-Go): the oracle ceiling is +0.08 on LongBench and +0.00 on RULER, and the +0.08 is fully explained by positive noise-selection bias (the median $|\Delta|$ over the 12 non-musique tasks is 0.17). The signal is available at zero cost, but the flat parameter surface leaves no leverage for adaptation. The best use of the same signal is the gate (safety value), not $\beta$ dispatch (no accuracy value). Refining the granularity to layers adds no leverage either (per-layer oracle gain averages +0.0008), and refining to the configuration-arm granularity reconfirms the rejection (13 end-to-end configuration arms): the per-arm oracle ceiling is $+0.21$ (versus fixed $\beta{=}0.25$) and $+0.49$ (versus the deployment-reference arm), none of the 12 candidate proxy features (task-matched mass, quality dispersion, synthetic-side mass, and so on) is significant after BH correction, the rank correlation between total mass and e2e accuracy is 0.166 (bootstrap 95\% CI containing 0), and picking arms by mass loses 1.18 against the oracle -- proxy-based arm selection is actively harmful, not merely neutral. The rejection at four granularities -- task, layer, sample, and configuration arm -- converges on one conclusion: \textbf{flatness is a structural property of the parameter surface, and no predictor has leverage to exert}, which translates into the deployment-simplicity argument.

\textbf{Dual-metric parameter grid and configuration election.} The election methodology of the main-table configuration: the offline quality grid scans a $9\times9\times9$ $(\alpha,\beta,\gamma)$ cube (0.125 step, 3600 valid configurations after legality-constraint filtering, 16-sample trace means, block pool $BP{=}64$, total budget $B_{\mathrm{TOK}}{=}2048$) over five method combinations (far/near scoring pairs (minmax, avg), (minmax, minmax), (avg, avg), (cluster, avg), (cluster, cluster)); the e2e grid curates 13 configuration arms, deduplicated by effective e2e budget from each combination's quality-optimal configurations, measured on LongBench hotpotqa and musique (n=200 each, $K_1{=}128$/$K_2{=}1024$). Two protocol lessons are registered honestly. First, the truncation collapse of $\gamma$: the 2$\times$ budget gap between the quality grid's $B_{\mathrm{TOK}}{=}2048$ and the e2e $K_2{=}1024$ makes $\gamma$ act on the effective budget through the $\min(\cdot,\,K_2^{\mathrm{mid}})$ truncation of the near fine-selection quota -- $\gamma$ arms that differ under the quality protocol (0.375 and 0.75) score bit-for-bit identically at e2e (54.16/35.57), so cross-protocol arm selection must deduplicate by the effective configuration of the target protocol. Second, the ranking reversal: the quality-protocol combination ranking (minmax, minmax) 0.896 $>$ (minmax, avg) 0.882 $>$ (cluster, avg) 0.879 $>$ (avg, avg) 0.796 reverses at e2e into (minmax, avg) 45.08 $>$ (minmax, minmax) 43.90 $\approx$ (avg, avg) 43.90 $>$ (cluster, avg) 43.58 (Figure~\ref{fig:reversal}); the (cluster, cluster) combination has no e2e implementation path on its near side (near-side clustering is unimplemented) and carries quality data only, and is excluded as such. The (cluster, avg) combination is the unified-pipeline counterpart of the ClusterKV-style cluster-representative flow (prefill semantic clustering with per-cluster selection~\cite{clusterkv}): its mass score of 0.879 nearly matches the (minmax, avg) 0.882, yet its e2e average trails by 1.50 (43.58 versus 45.08) -- the candidate-quality edge of semantic organization does not cash out under a tight token budget, consistent with the cluster negative-result verdict in \S4.4. Therefore \textbf{the election trusts only e2e measurements}: the winning arm (minmax, avg) with $\alpha{=}0.125/\beta{=}0.375/\gamma{=}0.625$ has the best average (45.08) and the highest hotpotqa of all 13 arms (55.44), with musique 34.71 second only to the $\beta{=}0.25$ arm's 35.57 (that arm averages 44.86, lower); it sits at $+0.49$ over the deployment-reference arm ($\gamma{=}0.125$) and $+0.21$ over the fixed $\beta{=}0.25$ -- the small magnitude of the gain itself reconfirms the flatness of the parameter surface. The complete configuration and per-task scores of the 13 arms are listed in Table~\ref{tab:e2egrid}: the selection rule takes the top-3 quality-optimal configurations per combination after deduplication by the effective e2e budget $(\alpha,\beta,\mathrm{nt\_near})$, plus the deployment-reference arm; the $\gamma$ truncation-collapse pair (0.375 and 0.75) is directly visible in the table as bit-for-bit identical scores.

\begin{table}[h]\centering
\small\setlength{\tabcolsep}{4pt}
\caption{Full list of the 13 e2e-grid configuration arms (LongBench hotpotqa/musique, n=200 each, $K_1{=}128$/$K_2{=}1024$; mass is the offline quality protocol, 16-sample trace means, listed for cross-reading but not used for election; the selection rule takes the top-3 quality-optimal configurations per combination after deduplication by $(\alpha,\beta,\mathrm{nt\_near})$ plus the deployment-reference arm; the mavg $\beta{=}0.25$ arms with $\gamma{=}0.375$ and 0.75 score bit-for-bit identically, the direct evidence of the $\gamma$ truncation collapse; the cavg combination is the unified-pipeline counterpart of the ClusterKV-style cluster representative).}
\label{tab:e2egrid}
\begin{tabular}{lllllll}
\toprule
combo & $\alpha$ & $\beta$ & $\gamma$ & mass & hotpotqa & musique \\
\midrule
mavg & 0.125 & 0.25 & 0.75 & 0.882 & 54.16 & 35.57 \\
mavg & 0.125 & 0.25 & 0.375 & 0.881 & 54.16 & 35.57 \\
mavg & 0.125 & 0.375 & 0.625 & 0.880 & \textbf{55.44} & 34.71 \\
mavg (ref.) & 0.125 & 0.375 & 0.125 & 0.873 & 54.43 & 34.76 \\
mminmax & 0.875 & 0.875 & 0.5 & 0.896 & 54.40 & 33.40 \\
mminmax & 0.75 & 0.75 & 0.5 & 0.895 & 54.38 & 33.24 \\
mminmax & 0.625 & 0.625 & 0.5 & 0.895 & 53.88 & 33.04 \\
aavg & 0.875 & 0.875 & 0.375 & 0.796 & 54.40 & 33.40 \\
aavg & 0.875 & 0.75 & 0.5 & 0.796 & 54.13 & 33.05 \\
aavg & 0.75 & 0.75 & 0.5 & 0.794 & 54.38 & 33.24 \\
cavg & 0.125 & 0.375 & 0.625 & 0.879 & 53.43 & 32.20 \\
cavg & 0.125 & 0.25 & 0.75 & 0.879 & 54.13 & 33.03 \\
cavg & 0.125 & 0.25 & 0.375 & 0.877 & 54.13 & 32.50 \\
\bottomrule
\end{tabular}
\end{table}

\begin{figure}[h]\centering
\includegraphics[width=0.95\linewidth]{figures/fig11_mass_vs_e2e.pdf}
\caption{Best configuration per combination under the quality-proxy protocol versus the e2e protocol. The dual-protocol difference must be stated: mass is softmax quality coverage (block pool $BP{=}64$, total budget $B_{\mathrm{TOK}}{=}2048$, 16-sample trace mean); e2e is the equal-weight mean of LongBench hotpotqa and musique (n=200 each) at $K_1{=}128$/$K_2{=}1024$. Both panels share the same $x$ order (sorted by mass descending), and the rank labels display the reversal directly: the mass champion (minmax, minmax) (0.896) drops to the third e2e tier (43.90, tied with (avg, avg)), while (minmax, avg), second under mass, becomes the e2e champion (45.08; highest hotpotqa of all 13 arms at 55.44, musique 34.71 second only to the $\beta{=}0.25$ arm's 35.57, which averages 44.86). White text inside the bars gives the two-task score of each combination's best e2e arm; the gray dashed line is the deployment-reference arm ($\gamma{=}0.125$, 44.59); the second row of the $x$ axis gives the best arm's $(\alpha,\beta,\gamma)$ under that protocol; (cluster, cluster) has no e2e implementation path and mass data only, excluded as such.}
\label{fig:reversal}
\end{figure}

\textbf{The offline-replay versus e2e protocol gap.} We double-ranked the five near/far scoring combinations under offline trace replay and end-to-end evaluation: the combination ranked 4th in replay quality (near-side mean aggregation, (minmax, avg), under the five-combination screening replay) overtakes to the top at e2e -- the rankings of the two protocols reverse systematically. The reversal is then quantified on the 13-arm dual-metric grid (Figure~\ref{fig:reversal}): under the grid replay protocol the same combination ranks 2nd of the four (a different rank than under the screening replay, with the reversal itself unchanged), and the rank correlation between the quality protocol and the e2e protocol is only 0.166 (bootstrap 95\% CI containing 0), the quality-protocol champion drops to the third e2e tier, and picking arms by the quality protocol loses 1.18 against the e2e optimum -- the proxy protocol is not merely unstable in ranking usable-level combinations, it is actively misleading. The gap reappears along the subspace-selection axis: the per-layer static-pair arm gains $+3.8$ points in replay ($R_\mathrm{far}$ 0.849 versus 0.811) yet loses $-0.50$ at e2e in the same configuration (50.04 versus 50.54, musique $-5.27$) -- selection bases derived from prefill statistics (trained projections and static pairs are structurally homologous) fail under decode query drift. This shows that offline proxy metrics cannot replace end-to-end evaluation, and that kernel microbenchmarks and e2e must both be measured and both reported. The regularity of the gap shows up in two further pieces of evidence. First, the trace mass of the top-$\sigma$ arms rises with $\sigma$ (0.7352 $\rightarrow$ 0.7522) while e2e falls monotonically -- the replay protocol gets the direction of marginal far-retrieval-quality change wrong. Second, among the five subspace e2e arms, rope64 ranks last in the trace 16-dimensional protocol yet sets the new e2e high. The reliability boundary of the two protocols follows: trace mass reliably detects collapse-level differences (random and highfreq collapse under both protocols) but cannot rank usable-level combinations -- ablation screening uses trace to exclude collapse, and the final verdict must be e2e.

\subsection{Kernel-Level Microbenchmarks}

The selection chain achieves the following speedups over dense full-dimensional scoring: 1.46$\times$ at 32K, 2.77$\times$ at 64K, and 5.09$\times$ at 128K (synthetic data, labeled as such, trace-replay protocol). The kernel and storage numbers of this subsection correspond to the production configuration (tail32-4bit L1), whose 13-task quality is 50.53 ($+0.16$ against FullKV with a CI containing 0, parity with the full-dimension L1 main-table arm; see the dual-protocol verification in \S4.2). The same-machine three-way comparison (official kernels imported as-is into a unified harness, 131K with cudaEvent timing) gives Quest 0.107ms, DSA 0.503ms, and PSI 0.787ms. The latency ranking is Quest $<$ DSA $<$ PSI, which we report as measured; the current advantage of PSI sits on the algorithm-structure side: $\approx$258 MACs per token for indexing (1/32 of DSA), index storage 3$\times$ below Quest (336B/token versus about 1KB/token, total over heads: Quest indexes all 32 heads while PSI indexes only the 8 kv-heads), and quality versus Quest at $+3.06$ on LongBench (main-arm protocol) and $+7.98$ on RULER (reference-arm protocol). At 131K all three sit far from the HBM bound, so the ranking reflects implementation maturity, not the algorithmic ceiling.

\subsection{End-to-End Speed}

The e2e three-arm verdict (Qwen3-30B-A3B, 2$\times$H20, same machine, same protocol, same env: PSI versus the Quest backend versus dense Triton, with prefill/total/decode measured separately by the two-pass method): Quest wins e2e in every shape -- a single 64K request totals 14.9s (dense 23.6s, PSI 33.3s), and TP2 bs16 64K totals 62.1s (dense 106.1s, PSI 142.0s). The loss of PSI concentrates in prefill (TP2 136.0s versus 54.4s for Quest: the per-chunk construction cost of the two-level index, while Quest builds only the page-level min/max single-level index); Quest's prefill is in turn faster than dense Triton (54.4s versus 102.3s), which shows that the three prefill paths are not compute-equivalent, so the prefill segment is read as a backend-to-backend comparison, and the prefill attribution of PSI (two-level index construction) holds within its own path. The decode segment reverses the ranking -- dense 61.6 $<$ PSI 93.9 $<$ Quest 121.1 ms/step; the relative advantage of PSI in the decode segment exists only in the increment protocol against Quest (index overhead 32 versus 59.5ms/step, 1.85$\times$) while it still trails dense by 32.3ms/step -- the decode segment is the only shape where PSI holds an advantage over Quest, not a net-win shape; in line with the selection-chain kernel speedup (1.46$\times$ at 32K to 5.09$\times$ at 128K), and combined with the validated cross-request batching (3.8$\times$) it is the path toward a net win. Two fix paths toward a net e2e win are identified: cross-request batching (validated at 3.8$\times$) and kernelizing the prefill index construction. The 8B steady state (bs16 $\times$ S=30K $\times$ n=256; the 733.9s and 65.0ms baselines are PSI itself before kernelization, not dense): decode step 65.0 $\rightarrow$ 33.5ms (1.94$\times$ self-improvement, 1.21$\times$ over the dense 40.4ms); prefill 733.9 $\rightarrow$ 183.4s (4.00$\times$ self-improvement, still 1.60$\times$ slower than dense, because the 8 kv-heads of the 8B model make the per-layer gather heavier; under the same protocol the 30B TP2 4-head shape takes 1.33$\times$ the dense prefill time (33\% slower), so the kv-head count narrows but does not reverse the prefill gap, and prefill index construction is the current dominant bottleneck).

\subsection{Profiling Methodology Findings}

Per-kernel microbenchmarks and production e2e differ systematically: the three host synchronizations on the select hot path (.item(), .any()) cause $\sim$100K queue drains in production form. A synthetic bench without a GPU queue measures none of this, and removing the synchronizations improves prefill by 6.7\% on the same workload. Synchronization-elimination optimizations must be validated at the e2e layer. We combine this with the observation that offline trace-replay rankings cannot substitute for e2e verdicts into one methodology: measure and honestly report both the kernel level and the e2e level.

% =====================================================================
% 4 Related Work: three-way clustering + positioning sentence (MoBA style:
% Category -> Limitation -> Position).
% =====================================================================
\section{Related Work}

\textbf{Sparse attention indexing.} Quest~\cite{quest} constructs a lower-bound approximation from page sampling, and the KV compression family (SnapKV, H2O, PyramidKV~\cite{snapkv,h2o,pyramidkv}) discards KV once at prefill. PSI differs from them on three counts: the score representation takes an upper bound (versus a lower bound: the upper-bound family leads Quest on the overall RULER average, by $+8.40$ under the TIA protocol and $+7.98$ under the PSI reference arm, the dominant source of the long-context quality split), the budget organization takes a partition (versus a single pool), and the structure is a two-level cascade (versus a single level). While adapting SnapKV, H2O, and PyramidKV under a unified protocol, we found that implementations without forced sink retention emit EOS as their first token on Qwen3-8B: sink receives no selection signal in shallow layers yet carries 0.59 of attention mass in deep layers, a layer-by-layer progressive distortion. This is both a fairness correction for adapting this family of baselines and corroborating evidence for forced sink retention.

\textbf{Training-based indexing.} DSA~\cite{dsa} (DeepSeek V3.2) trains a selector with $I=\sum w\cdot\mathrm{ReLU}(q\cdot k)$, 64 heads $\times$ 128 dimensions in FP8 with top-2048, and a 1000-step warm-up; MoBA~\cite{moba} applies gated block sparsity. PSI is training-free and weight-free, and we compare against them in a unified kernel-level harness. The two form a methodological divide: DSA trades training for selection capability and can learn arbitrary projections, while a training-free method can rely only on position priors that are independent of the query distribution, and the protocol gap (the trace protocol holds while e2e collapses) is the empirical characterization of this dividing line.

\textbf{Two-level structures.} HISA~\cite{hisa} (COLM 2026) validates the block-coarse-plus-token-refine form; InfLLM~\cite{infllm} builds a context memory of block summaries for two-level retrieval, Landmark Attention~\cite{landmark} achieves random-access context through landmark block representatives, and ClusterKV~\cite{clusterkv} organizes the KV cache as prefill semantic clusters selected per cluster -- so the two-level / block-unit / cluster-unit selection form itself is not our contribution; for the cluster representatives the cluster center is not an upper bound and carries a theoretical miss hole. The increment of PSI is the systematic characterization of two orthogonal dimensions: representation (upper bound versus lower bound -- none of the above carries an upper-bound guarantee) and budget organization (partition versus single pool), and their quality-speed impact.

\textbf{Inference systems.} The FlashAttention~\cite{fa} line optimizes compute kernels, and PagedAttention/vLLM~\cite{vllm} optimizes memory organization. PSI is orthogonal to both and integrates into sglang in production form.

% =====================================================================
% 5 Limitations: honest boundaries.
% =====================================================================
\section{Limitations and Future Work}

On accuracy, the main-table LongBench arm differs from FullKV by $+0.42$ (8 wins and 5 losses over 13 tasks, all losses $\leq$1 point) -- the 13 tasks carry no seed repetition, and the 95\% per-sample paired-bootstrap CI of the $+0.42$ difference is $[-0.18, +1.02]$, containing 0 (held-out 11 tasks: $+0.08$, $[-0.46, +0.60]$), so the accuracy claim is parity-level rather than a significant lead; the best arms of the two benchmarks belong to two different $(\beta,\gamma)$ configurations (the LongBench real-task profile at $\beta{=}0.375/\gamma{=}0.625$, the RULER synthetic-needle profile at $\beta{=}0.25/\gamma{=}0.125$) -- the same-table RULER rescan of the main arm (a full 33-arm sweep) lands at 85.83 ($-2.10$; the six-arm 32K supplementary point extrapolates it to $-2.15$, stable with length), with the damage concentrated in the far-dense word-list task cwe (82.0/68.6/55.5, absolute score degrading over 4K--16K and recovering to 61.9 at 32K) and 7 of the other 9 tasks matching exactly with multikey\_3/multivalue within 3 points: the uniformity question is now verified and reported as measured, and the deployment protocol is to select one of the two arms by task profile; the best RULER arm at 87.93 still trails FullKV at 88.71 ($-0.78$) and the first-generation TIA at 88.35 ($-0.42$), with the residual cost concentrated in multikey\_3 (16K: 75--80 versus 89 for TIA), where the 4bit budget granularity dominates on tasks with an extremely large number of distractor keys. On speed, the e2e three-arm comparison leaves PSI behind Quest and dense in every shape (TP2 142.0s versus 62.1s and 106.1s), with the loss concentrated in the prefill two-level index construction (136.0s versus 54.4s for the Quest single-level index); the 1.85$\times$ index-overhead advantage over Quest in the decode segment is the sole e2e cash-out point; cross-request batching (validated at 3.8$\times$) and kernelizing the prefill index construction are the two fix paths. On scale and model family, all conclusions are obtained on the single Qwen3 family -- the subspace mechanism weakens on the 32B member of the same family (entry-recall protocol: 0.66--0.87 with a gap of $\le$0.08 to the full dimension; the mechanism exists, reported as measured), but cross-family extrapolation is unverified: the rotation-pair mechanism depends on the tail-dimensions-equal-complete-low-frequency-pairs correspondence of the rotate\_half (GPT-NeoX-style) layout, which does not hold under interleaved (GPT-J-style) layouts and requires rederiving the dimension selection, and the position-stable low-frequency subspace premise itself is untested on RoPE-free or partial-rotation models. Latent indexing for MLA architectures (the DeepSeek family) is future work. The no-miss upper-bound property holds only within the subspace-scoring protocol; cross-protocol misses have static-pair-arm-type evidence (replay gains do not extrapolate). The far statistic of the gate is a single per-layer value, which cross-request contamination affects in mixed batches (production semantics should be per-row).

% =====================================================================
% Appendices.
% =====================================================================
\section*{Appendix A: Index of Experiment Identifiers}
The mapping from scripts to conclusions for all experiment identifiers and system milestones (including commit hashes, cross-check protocols, and registered negative results) is maintained in the companion experiment log shipped with the source code.

\section*{Appendix B: Reproducibility}
2$\times$H20-3e (141GB each), torch 2.8.0+cu128, sglang two-level-indexer branch; RULER official pre-generated set (KVCache-Factory mirror), LongBench v1 full, and the real weights of Qwen3-8B, 30B-A3B, and 32B; string\_match\_all official scoring; e2e differencing by alternating sparse and dense configurations in matched runs.


\begin{thebibliography}{99}\small

\bibitem{dsr1}
DeepSeek-AI. DeepSeek-R1: Incentivizing Reasoning Capability in LLMs via Reinforcement Learning. arXiv:2501.14348, 2025.

\bibitem{qwen3}
Qwen Team. Qwen3 Technical Report. arXiv:2505.09388, 2025.

\bibitem{fa}
T.~Dao, D.~Y. Fu, S.~Ermon, A.~Rudra, and C.~R\'{e}. FlashAttention: Fast and Memory-Efficient Exact Attention with IO-Awareness. In \emph{NeurIPS}, 2022.

\bibitem{vllm}
W.~Kwon, Z.~Li, S.~Zhuang, Y.~Sheng, L.~Zheng, C.~H. Yu, J.~E. Gonzalez, H.~Zhang, and I.~Stoica. Efficient Memory Management for Large Language Model Serving with PagedAttention. In \emph{SOSP}, 2023.

\bibitem{quest}
Y.~Tang, L.~Lin, Y.~Cao, S.~Li, F.~Yang, Y.~Wang, and T.~Wu. Quest: Query-Aware Sparsity for Efficient Long-Context LLM Inference. In \emph{ICLR}, 2025.

\bibitem{snapkv}
Y.~Li, Y.~Huang, B.~Yang, Y.~Li, Y.~Liu, and M.~Qiu. SnapKV: LLM Knows What You Are Looking For Before Generation. In \emph{NeurIPS}, 2024.

\bibitem{h2o}
Z.~Zhang, Y.~Sheng, T.~Zhou, D.~Chen, L.~Zheng, B.~Chen, L.~Cai, M.~Yan, Y.~Yu, J.~Tang, J.~Xu, P.~Zhang, and R.~Chen. H2O: Heavy-Hitter Oracle for Efficient Generative Inference of Large Language Models. In \emph{NeurIPS}, 2023.

\bibitem{pyramidkv}
Z.~Cai, Y.~Zhang, B.~Wang, X.~Liu, H.~Zhao, Y.~Lu, Y.~Wang, X.~Sun, K.~Liu, and Z.~Duan. PyramidKV: Dynamic KV Cache Compression Based on Pyramidal Information Funneling. arXiv:2406.02069, 2024.

\bibitem{streamingllm}
G.~Xiao, Y.~Tian, B.~Chen, S.~Han, and M.~Lewis. Efficient Streaming Language Models with Attention Sinks. In \emph{ICLR}, 2024.

\bibitem{longformer}
I.~Beltagy, M.~E. Peters, and A.~Cohan. Longformer: The Long-Document Transformer. arXiv:1904.07785, 2020.

\bibitem{bigbird}
M.~Zaheer, M.~Guruganesh, K.~A. Dubey, J.~Ainslie, C.~Alberti, S.~Onta\~{n}\'{o}n, P.~Pham, A.~Ravula, Q.~Wang, L.~Yang, and A.~Ahmed. Big Bird: Transformers for Longer Sequences. In \emph{NeurIPS}, 2020.

\bibitem{sparq}
SparQ Attention: Accelerating LLM Inference with Accurate KV Retrieval, Minimized Decoding Costs, and No Fine-Tuning. arXiv:2406.09050, 2024.

\bibitem{dsa}
DeepSeek-AI. DeepSeek-V3.2-Exp Technical Report: Boosting Long-Context Efficiency with Exact Sparse Attention. 2025.

\bibitem{moba}
Moonshot AI. MoBA: Mixture of Block Attention for Efficient Long-Context LLMs. arXiv:2502.13189, 2025.

\bibitem{hisa}
Y.~Xu, F.~Meng, F.~Jiang, Y.~Wang, R.~Zhou, Z.~Wang, J.~Wu, Z.~Pan, X.~Tang, W.~Pei, T.~Liu, D.~Yin, X.~Sun, and M.~Zhang. HISA: Efficient Hierarchical Indexing for Fine-Grained Sparse Attention. In \emph{COLM}, 2026. arXiv:2603.28458.

\bibitem{infllm}
% TODO: verify the formal entry before submission
M.~Fan, W.~Wang, R.~Ren, X.~Lin, J.~Guo, F.~Wang, and T.-S.~Chua. InfLLM: Training-Free Long-Context Extrapolation for LLMs with an Efficient Context Memory. arXiv:2408.12819, 2024.

\bibitem{landmark}
% TODO: verify the formal entry before submission
B.~Chiang, S.~T. T. Le, and Moonton. Landmark Attention: Random-Access Infinite Context Length for Transformers. In \emph{ICLR}, 2024.

\bibitem{clusterkv}
% TODO: verify the formal entry before submission (paper page unreachable in this environment; arXiv ID and authors to be filled)
ClusterKV: Manipulating LLM KV Cache in Semantic Space. 2025.

\bibitem{ruler}
C.-P. Hsieh, S.~Sun, S.~Kriman, S.~Ganti, J.~Zhang, B.~Kulis, K.~Papayiannis, M.~Fofanos, N.~Rudrajit, and C.~Tigges. RULER: What's the Real Context Size of Your Long-Context Language Models? arXiv:2404.06654, 2024.

\bibitem{longbench}
Y.~Bai, X.~Lv, J.~Zhang, H.~Zhao, X.~Qi, L.~Huang, Z.~Du, X.~Wang, Y.~Zheng, Y.~Zheng, J.~Zhang, Y.~Li, T.~Su, and J.~Tang. LongBench: A Bilingual, Multitask Benchmark for Long Context Understanding. In \emph{ACL}, 2024.
\end{thebibliography}

\end{document}
